\documentclass[11pt,letterpaper]{article}
\usepackage[T1]{fontenc}
\usepackage{lmodern,microtype}
\usepackage[margin=1in]{geometry}
\usepackage{amsmath,amssymb,amsthm,mathtools}
\usepackage{booktabs,tabularx,longtable,array,enumitem,xcolor}
\usepackage{xurl,hyperref,fancyhdr,listings}
\definecolor{navy}{HTML}{183653}
\definecolor{teal}{HTML}{176B76}
\hypersetup{colorlinks=true,linkcolor=navy,urlcolor=teal,citecolor=teal,
 pdftitle={Proof Transport for the 215-State ZFC Machines},
 pdfauthor={ChatGPT research contribution for ProveIt}}
\pagestyle{fancy}
\fancyhf{}
\fancyhead[L]{\small\sffamily\color{navy}PROVEIT / COMPUTABILITY}
\fancyhead[R]{\small\sffamily K215 proof transport}
\fancyfoot[C]{\small\thepage}
\setlength{\headheight}{14pt}
\setlength{\parskip}{5pt}
\setlength{\parindent}{0pt}
\setlist{itemsep=3pt,topsep=5pt}
\newtheorem{theorem}{Theorem}[section]
\newtheorem{lemma}[theorem]{Lemma}
\newtheorem{proposition}[theorem]{Proposition}
\newtheorem{corollary}[theorem]{Corollary}
\theoremstyle{definition}
\newtheorem{definition}[theorem]{Definition}
\newcommand{\Halts}{\operatorname{Halts}}
\newcommand{\Con}{\operatorname{Con}}
\newcommand{\Refute}{\operatorname{Ref}}
\newcommand{\ZFC}{\mathrm{ZFC}}
\newcommand{\K}{\mathrm{K215}}
\newcommand{\CM}{\mathrm{CM215}}
\newcommand{\repo}{https://github.com/VladimirReshetnikov/ProveIt/blob/b6c7087477e51c017003bd729e712c8969733223}
\lstset{basicstyle=\small\ttfamily,breaklines=true,columns=fullflexible,
 keepspaces=true,frame=single,rulecolor=\color{black!15},
 backgroundcolor=\color{black!2},showstringspaces=false}

\begin{document}
\begin{titlepage}
\vspace*{0.5in}
{\sffamily\large\color{teal}RESEARCH AND VERIFICATION NOTE\par}
\vspace{0.3in}
{\sffamily\bfseries\Huge\color{navy}Proof Transport for the\par
215-State ZFC Machines\par}
\vspace{0.25in}
{\Large A local CM215 certificate, a source-proof audit,\par
and bounded searches for 214 states\par}
\vspace{0.4in}
{\large Prepared for the ProveIt repository\par}
{\large 21 October 2026\par}
\vspace{0.35in}
\begin{abstract}
The most useful simplification for a second verified 215-state machine is
to transfer correctness directly from the already calibrated K215 table.
The existing conditional-marker candidate CM215 changes six transition
rows and keeps the source program unchanged. We replace its history-based
explanation by a three-clause inductive tape invariant. Each old transition
then has an exact one- or three-step implementation, preserving the entire
tape and head at the comparison points. A reusable positive-block theorem
turns this single simulation into halting equivalence in both directions.

The accompanying artifacts include the generic theorem checked in Lean
4.32.0 and Coq 8.18.0, the concrete Lean CM215 development, literal machine
tables, finite checks, and an interface to the repository's dense-machine
semantics. They separate new checks from the retained K215 calibration.
The current repository already incorporates three proof-simplification
phases; this report credits those changes and identifies the remaining
source-certificate opportunity: 2,240 zero-offset checkpoints use only 32
local arithmetic-effect signatures, and their labelled graph has a checked
1,423-class quotient.

No verified 214-state machine was obtained. Completed searches evaluate
17,708,928 fold specifications across fourteen literal tables and two
initial domains, plus 159,509,360 individual branch specifications in a
second restricted family. None survives the stated finite-window filters.
These are reproducible search boundaries, not state lower bounds.
\end{abstract}
\vfill
\textbf{Pinned repository snapshot}\par
\small\texttt{b6c7087477e51c017003bd729e712c8969733223}\par
\normalsize This contribution concerns proof reuse and verified local
transformation. The 215-state tables themselves and their foundational
calibration are inherited from ProveIt and its credited upstream authors.
\end{titlepage}

\tableofcontents
\newpage

\section{Answers and evidence boundaries}

The recommended next step is to retain the checked K215 calibration as a
base theorem and attach small machine-to-machine proofs for local variants.
For CM215 this removes the need to redo the source enumerator, numerical
calculus, register compiler, 4,751 primitive-source occurrences, and blank
initialization. It reduces the \emph{additional proof} needed for another
candidate. The base mathematics remains in the final dependency chain.

The question about fewer states has a different outcome: this work does
not supply a correct 214-state machine. It supplies complete, bounded,
reproducible searches and a reusable way to certify a future successful
local transformation.

\begin{center}
\begin{tabularx}{\textwidth}{@{}p{0.27\textwidth}X@{}}
\toprule
Result & Status and scope\\
\midrule
Exact K215 consistency endpoint & Existing repository theorem and retained
Lean/Rocq build evidence; the complete baseline was not rebuilt here.\\
Positive-block transport & New small checked artifact in Lean 4.32.0 and
Coq 8.18.0; the Coq object also passed recursive \texttt{coqchk}.\\
CM215 local transformation & Three-clause invariant, exact whole-tape
blocks, and a concrete Lean transfer development. The literal candidate
already existed in the repository.\\
Source-certificate analysis & Fresh finite computation: 2,240 entries,
zero nonzero offsets, 32 effect signatures, 1,423 quotient classes.
This is not yet a replacement kernel checker.\\
Search below 215 states & Two precisely limited families completed;
no 214-state candidate obtained and no minimality claim.\\
\bottomrule
\end{tabularx}
\end{center}

\subsection{The target statement}

Let $\Refute_{\ZFC}$ mean that there is a finite syntactic derivation of
falsity in the repository's ZFC presentation, and let
$\Con(\ZFC)\equiv\neg\Refute_{\ZFC}$. The base endpoint states
\[
 \Halts(\K,\mathrm{blank})\ \longleftrightarrow\ \Refute_{\ZFC}.
\]
No consistency hypothesis is a premise of this equivalence. Combining it
with a proof that K215 and CM215 halt equivalently gives
\[
 \neg\Halts(\CM,\mathrm{blank})\ \longleftrightarrow\ \Con(\ZFC).
\]
The repository records a distinction between an external proof-assistant
theorem and a derivation encoded in ZFC's own proof syntax. We preserve it.
The new finite-machine argument can be formalized in the same elementary
arithmetic setting as the original machine argument; this packet does not
claim to have constructed a new encoded internal ZFC or PA derivation.
The existing ordinary ZFC-provability conclusion transports by elementary
propositional reasoning. See Section~\ref{sec:internal}.

\subsection{Machine model and identities}

All machine counts use one two-way infinite tape, two symbols, blank
symbol zero, and only left/right moves. Initialization is part of the
working-state count; the distinct halt is excluded. Both dense tables have
working states $0,\ldots,214$, initial control zero, and halt 215. Their
native tables preserve historical labels and use halt 254.

The dense K215 table has SHA-256
\begin{lstlisting}
fdbcef0825886035c3091c8537a69535bc96f9e265f0ae45863543d40ae70090
\end{lstlisting}
and the dense CM215 table has SHA-256
\begin{lstlisting}
1f4f084f7293e2b042d5746039f793ff72ed429083b504970c939a5515e666cf
\end{lstlisting}
Both files are included, with 430 rows each. The root manifest binds the
literal data, new proof sources, and all retained evidence. Machine hashes
identify the theorem's target; a same-size table is not interchangeable
with a calibrated one.

\section{What the current proofs already achieve}

The pinned snapshot already contains three simplification phases. The
verification index records an original fresh Lean build of 204 project
modules and two explicit clients, and a Rocq build of 237 selected sources
with recursive checking. The current Lean closure is 185 project modules.
The phase-three record rebuilt nine affected modules and reused 176
hash-bound objects; it is an incremental check, not a second full replay.
The final Lean reports use only \texttt{propext},
\texttt{Classical.choice}, and \texttt{Quot.sound}; the four retained
Rocq endpoint dependency reports are closed under the global
context.\cite{verification,phase3}

\begin{center}
\begin{tabular}{@{}lrr@{}}
\toprule
Proof source & Original bytes & Current bytes\\
\midrule
Lean dominated implication & 689,857 & 35,185\\
Lean AX routine & 202,222 & 7,563\\
Lean source-certificate data & 653,569 & 456,548\\
Lean site-coverage facade & 107,243 & 14,916\\
Lean register determinism & 17,209 & 12,988\\
Rocq dominated implication & 921,691 & 22,673\\
\bottomrule
\end{tabular}
\end{center}

These are repository-recorded source-size reductions, not performance
improvements measured in this session. In particular, symmetry already
removed the need for sixteen reverse source-certificate check modules.
The latest Lean dominated-rule development already uses 61 bootstrap
lemmas, 81 conditional schemas, and a deduction theorem. The existing
enumerator proof already shares one theorem between the zero and boot
starts. The site facade already directly composes 105 checked blocks.
Recommending these as new work would duplicate completed work.

\subsection{The endpoint facade is not the expensive part}

The Lean endpoint file is 53 lines. The corresponding Rocq sparse and
dense endpoint facades are 19 and 18 lines. The Lean facade supplies the
following machine-specific inputs to a shared theorem:

\begin{enumerate}
\item A commuting dense-state renaming into the sparse machine table.
\item Agreement with already proved register-engine contracts.
\item A blank bootstrap reaching the correct dispatcher frame.
\item Coverage of every literal source occurrence.
\item Source/program correspondence and the numerical-calculus
      refutation equivalence.
\end{enumerate}

The final composition is already appropriately small. The profitable
targets are generated certificates, repetitive local contracts, and
repeated verification of closely related candidates.\cite{endpoint}

\subsection{Proof reuse should follow the kind of change}

There are three different situations. A local table rewrite with unchanged
source is best handled by direct machine simulation. A source reassociation
with the same command semantics is best handled by a source control-flow
certificate plus existing engine contracts. A new command language or
proof calculus needs a semantic calibration at that higher level before
the compiler can be reused.

This separation prevents a small physical edit from triggering a new proof
of set theory. It also prevents an apparently shorter source from being
accepted merely because its table and startup look plausible. The proof
interface must expose every changed semantic assumption.

\section{A reusable positive-block theorem}\label{sec:positive}

Let $f:A\to A\cup\{\bot\}$ and $g:B\to B\cup\{\bot\}$ be deterministic
partial transition functions. The value $\bot$ means that no transition
exists. Write $f^n(a)\downarrow a'$ for a successful $n$-step execution;
this includes the assertion that no earlier step became undefined. Define
\[
 \Halts_f(a)\quad\Longleftrightarrow\quad
 \exists n,a'\;(f^n(a)\downarrow a'\ \land\ f(a')=\bot).
\]

\begin{theorem}[Positive-block transport]\label{thm:transport}
Suppose a relation $R\subseteq A\times B$ has the following properties
whenever $a\,R\,b$:
\begin{enumerate}
\item $f(a)=\bot$ if and only if $g(b)=\bot$.
\item If $f(a)=a'$, then for some $k>0$ and $b'$,
      $g^k(b)\downarrow b'$ and $a'\,R\,b'$.
\end{enumerate}
Then $a\,R\,b$ implies $\Halts_f(a)\leftrightarrow\Halts_g(b)$.
One forward simulation suffices; no reverse certificate is needed.
\end{theorem}

\begin{proof}
For preservation of halting, induct on the number of source steps before
the terminal configuration. At zero steps use terminal agreement. At a
positive step, the second condition provides a successful target block.
Concatenate that block with the terminating target execution provided by
the induction hypothesis at its endpoint.

For reflection, fix a terminating target run of length $n$ and use strong
induction on $n$. If the source is terminal, it already halts. Otherwise
the second condition supplies a successful target prefix of length $k>0$.
Determinism forces it to agree with the corresponding initial part of the
chosen terminating run. It cannot pass the terminal endpoint, because a
terminal configuration has no successor. Thus $n=k+m$ for some $m<n$,
and the target terminates from $b'$ after $m$ further steps. Apply the
induction hypothesis to $a'\,R\,b'$ and then prepend the one source step.
\end{proof}

\subsection{Why positivity matters}

If zero-length implementations were allowed, let the source loop at $a$
forever, and let the target take one transition from $b$ to a terminal
state. Relate only $a$ and $b$. Terminal agreement holds at the related
pair, and every source step is represented by zero target steps, yet only
the target halts. A well-founded auxiliary ranking could replace strict
positivity in a more general theorem; it cannot simply be omitted.

There is no fairness assumption in Theorem~\ref{thm:transport}. The result
concerns deterministic partial systems, and successful positive prefixes
provide exactly the progress required by the reverse induction. There is
also no uniform upper bound on $k$ in the theorem. CM215 happens to use
only $k=1$ and $k=3$.

\subsection{Checked implementations}

\path{formal/PositiveStutter.lean} imports only \texttt{Init} and implements
the theorem for arbitrary configuration types. Its halting and nonhalting
equivalences were compiled with Lean 4.32.0. Their printed axiom
dependencies are \texttt{propext} and \texttt{Quot.sound}, with no
\texttt{sorryAx}.

\path{formal/PositiveStutter.v} gives the same theorem in the Coq/Rocq
language. It was compiled here with \emph{Coq 8.18.0}, and both public
equivalences print ``Closed under the global context.'' Its compiled
object passed \texttt{coqchk}. The pinned repository uses Rocq 9.1.0 for
its full baseline; this session's generic Coq check is stated with its
actual version and is not represented as a new full Rocq endpoint build.

The theorem is a standard simulation principle. The contribution is a
small reusable implementation, its application to the literal CM215
transformation, and the exact separation of local obligations from the
large existing calibration.

\section{CM215: a three-clause invariant}

CM215 is the existing conditional-marker alternative. It shares the
source program with K215, including its 202 distinct source subtrees and
4,751 primitive occurrences. It changes six native transition rows.
Its original packet supplies an ordinary whole-tape simulation theorem
and separate runtime evidence.\cite{cm215}

\begin{center}
\begin{tabular}{@{}lll@{}}
\toprule
Native state/read & K215 action & CM215 action\\
\midrule
$52/0$  & $0\;R\;154$ & $0\;L\;176$\\
$156/1$ & $1\;R\;176$ & $1\;R\;52$\\
$165/0$ & $0\;R\;176$ & $0\;R\;52$\\
$176/0$ & $0\;R\;168$ & $0\;R\;185$\\
$176/1$ & $1\;R\;150$ & $1\;R\;129$\\
$193/0$ & $0\;R\;176$ & $0\;R\;52$\\
\bottomrule
\end{tabular}
\end{center}

CM215 merges the old operational roles of 52 and 176 into control 52,
then reuses physical control 176 as a selector. This is a state-neutral
transformation: it leaves 215 working states. It does not itself save the
additional state sought by the user.

For a configuration $(q,h,t)$, let $t:\mathbb Z\to\{0,1\}$ be the entire
tape. There is no finite-support restriction in the local theorem.

\begin{definition}
The invariant $I(q,h,t)$ is the conjunction
\begin{align*}
 q=52  &\ \Longrightarrow\ t(h-1)=1,\\
 q=176 &\ \Longrightarrow\bigl(t(h)=1\ \lor\ t(h-1)=0\bigr),\\
 q=156 &\ \Longrightarrow\ t(h+1)=1.
\end{align*}
\end{definition}

These three clauses replace a history-sensitive explanation by a
one-step invariant. The third clause records exactly the information
needed for the exceptional incoming path into 176; no unbounded history
or auxiliary trace is retained.

\begin{lemma}[Invariant preservation]\label{lem:inv}
Every K215 execution starting in control zero satisfies $I$, for any
initial head position and any initial binary tape.
\end{lemma}

\begin{proof}
At control zero all three implications are vacuous. For preservation,
inspect the complete incoming rows of the three distinguished controls.
Control 52 has incoming rows $190/1$ and $193/1$, each writing one and
moving right. Hence the new cell immediately to the left is one.

Control 156 has the unique incoming row $189/1$, which writes one and
moves left. Hence the new cell immediately to the right is one.

Control 176 has incoming rows $165/0$, $193/0$, and $156/1$. The first
two write zero and move right, establishing the left-zero disjunct. The
last moves right from 156. The induction hypothesis for 156 says that
the cell it is about to read is one, establishing the current-one
disjunct. No other destination can make one of the implications
nonvacuous. Induction on the execution length completes the proof.
\end{proof}

\subsection{Exact local blocks}

Define $\phi(176)=52$ and $\phi(q)=q$ otherwise. At a comparison point
relate $(q,h,t)$ to $(\phi(q),h,t)$, with $I(q,h,t)$ required on the old
configuration.

For every old control outside $\{52,176\}$, the new row is the old row
with its target relabelled by $\phi$. Thus the corresponding one-step
executions agree on the whole tape and head. At both exceptional old
controls the read-one row is $1\,R\,150$, and the new control 52 has
that same row. These cases also take one step.

When old control 52 reads zero, the invariant gives $t(h-1)=1$, and the
new computation is
\[
 52\text{ at }h
 \ \xrightarrow{0L}\ 176\text{ at }h-1
 \ \xrightarrow{1R}\ 129\text{ at }h
 \ \xrightarrow{0R}\ 154\text{ at }h+1.
\]
Each written bit equals the bit already present. This is exactly the
head, tape, and destination of the old $52/0$ step.

When old control 176 reads zero, its invariant forces $t(h-1)=0$.
The new computation is
\[
 52\text{ at }h
 \ \xrightarrow{0L}\ 176\text{ at }h-1
 \ \xrightarrow{0R}\ 185\text{ at }h
 \ \xrightarrow{0R}\ 168\text{ at }h+1.
\]
Again the entire tape is unchanged and the endpoint is exactly the old
one-step endpoint. The auxiliary rows $129/0$ and $185/0$ are unchanged
rows of the literal candidate. The finite certificate checks them too.

\begin{theorem}[Whole-tape CM215 equivalence]\label{thm:cm}
For every initial binary tape $t$ and every head position $h$,
\[
 \Halts(\K,(0,h,t))\ \longleftrightarrow\
 \Halts(\CM,(0,h,t)).
\]
At corresponding comparison points the heads and complete tapes agree.
If $z(n)$ of the first $n$ old steps read zero in control 52 or 176, the
corresponding new prefix has exactly $n+2z(n)$ steps.
\end{theorem}

\begin{proof}
Use the relation just defined. Lemma~\ref{lem:inv} supplies its inductive
invariant. The row identities and two local blocks establish positive
advance with length one or three. The literal tables have matching
terminal behavior under $\phi$. Apply Theorem~\ref{thm:transport}.
The exact duration formula follows by adding the lengths of the local
blocks. In particular, $n\le n+2z(n)\le3n$.
\end{proof}

The theorem quantifies arbitrary initial tapes \emph{at control zero}.
It does not assert that the invariant holds from an arbitrary control
inside the modified fragment. For example, the comparison block for
old control 52 reading zero uses the left-one clause essentially.

\subsection{The concrete finite evidence}

The independent Python checker enumerates all 1,710 admissible three-bit
windows at the 215 working controls. Exactly four of these cases use
three new steps; all other cases use one. The four count both values of
the unused right neighbor in each of the two exceptional zero cases.
Terminal behavior is checked separately. The window enumeration verifies
the local hypotheses; the Lean induction extends them to arbitrary
infinite tapes and arbitrarily long runs.

The formal table facts range over all 216 dense control values and both
read bits, including the halt. Dense labels 52, 170, and 151 correspond
to native controls 52, 176, and 156. Kernel computations are partitioned
into small blocks; no \texttt{native\_decide} or unproved assertion is
used to accept the machine identities.

\subsection{Connection to the existing endpoint}\label{sec:internal}

The concrete development uses function-valued tapes because the local
invariant is particularly simple there. The repository's dense
\texttt{Machine 215} represents a finite-support tape by a list of
one-cell positions and uses \texttt{none} for the halt state. The bridge
maps that tape to its read function and maps \texttt{none} to dense label
215. It proves that read-after-write and head movement commute with the
embedding. The absorbed halt in the repository's total step is
partialized before applying the positive-block theorem.

The old literal table is compared with the exact
\texttt{KDense215.machine215}, not with a separately assumed equivalent
table. A companion dense machine is defined from the new literal rows.
This provides a direct blank-halting equivalence in the repository's
machine semantics. The small integration wrapper composes that
equivalence with the existing K215 ZFC theorem. The integration wrapper
itself was not compiled in this session. The freshly checked results are
the Core-machine equivalence and the parameterized
\texttt{calibration\_transfer} theorem. The complete 185-module calibration
was not rebuilt in this packet.

At the ordinary internal level, if ZFC proves
\[
 \Halts(\K,\mathrm{blank})\leftrightarrow
 \Halts(\CM,\mathrm{blank})
\quad\text{and}\quad
 \neg\Halts(\K,\mathrm{blank})\leftrightarrow\Con(\ZFC),
\]
then ZFC proves the requested CM215 equivalence by transitivity and
negation. The local tape proof uses finite transition inspection and
induction on natural-number execution lengths. It introduces no new
set-theoretic axiom and no additional consistency premise.

\section{Which 215-state candidate is easier to verify?}

For adding a second verified endpoint, CM215 has the clearest advantage:
its entire new burden is a local transport proof from a checked base.
This conclusion is about proof reuse, not merely source size.

\begin{center}
\begin{tabularx}{\textwidth}{@{}p{0.19\textwidth}p{0.30\textwidth}X@{}}
\toprule
Candidate & Relevant design change & Consequence for a new proof\\
\midrule
CM215 & Six rows; same source; local one/three-step simulation & Reuse
the complete K215 calibration, source proof, and startup. Best immediate
choice for an additional endpoint.\\
Rebalanced R41 & Nineteen native rows; 4,120 source addresses changed;
maximum depth 47 to 41 & Useful address-complexity reduction, but requires
a new source correspondence and runtime bindings.\\
Commands120 & DROP/NIP/MP command ordering & Reuses the same broad
calculus but changes source/control and startup evidence.\\
Direct-transfer ST215 & A direct 35-leaf transfer command and changed
source association & Fewer primitive occurrences, but an additional
command/compiler and protected-tail argument.\\
Zero-spare variants & Changed spare representation and startup seed &
Additional representation and startup obligations; not automatically
easier because the stored constants are smaller.\\
\bottomrule
\end{tabularx}
\end{center}

The original CM215 packet reports 6,176 runtime cases and 1,701,542
symbolic rules, slightly more than the primary machine's 1,698,790 rules.
Thus it would be misleading to justify CM215 as easier merely by its
full runtime-certificate size. Its advantage is that this second full
runtime proof can be bypassed by the direct equivalence theorem.
The primary K215 is still the established starting point for a complete
calibration from first principles.\cite{cm215,verification}

R41 is a sensible next subject if the aim is to simplify the decoder and
source checker themselves. The retained packet reports 1,640,317 symbolic
rules, about 3.44\% fewer than the primary figure, and the same 95,714-step
bootstrap. That arithmetic comparison does not demonstrate a faster
kernel build.\cite{rebalanced}

ST215 has 4,732 primitive occurrences and 6,151 guarded cases, with a
70,506-step blank prefix. Its command change means those smaller numbers
must be weighed against a different semantic proof. The repository's
intake notes distinguish delivered claims from checks rerun at intake;
this report does not upgrade those notes into new independent formal
verification.\cite{directtransfer,commands120,zerospare}

\section{A concrete next simplification of the source proof}

\subsection{Remove unused offset generality carefully}

The current K215 source proof compares the optimized path program with
\texttt{DagPerm} using an offset segment-bisimulation certificate. Each
entry records two checkpoint paths, an integer offset on register five,
and two successor entries. All 2,240 offsets in this particular
certificate are zero.\cite{sourceproof}

The corresponding semantic relation is simply equality of register
files. A specialized checker can avoid integer casts, offset propagation,
and offset-zero side conditions. This should be implemented as an
instance of a proved general interface, so future offset-using candidates
can retain the existing checker.

There is a concrete trap: the existing zero-offset checker decodes shorter
path codes than the offset checker. Our replay finds maximum checkpoint
path lengths 45 and 176 on the two sides; the target DAG has height 178.
The old zero-offset decoding bound is insufficient. Parameterize the
decoder depth, or use a structurally terminating path representation,
before switching this source proof to zero-offset checking. Dropping the
offset field without preserving path decoding would change the theorem.

\subsection{Separate local effects from addresses}

Between register tests, a segment contains only increments and control
movement. Its action on the register file is captured by an increment
count for each affected register. A checkpoint can therefore carry:
\[
 \bigl(\text{tested register},\ \mathbf v_{+},\ \mathbf v_{0},
        \text{positive successor},\ \text{zero successor}\bigr),
\]
or a halt label. The positive branch first performs the guarded decrement;
the vectors describe increments in the following segment. This is the
semantics already exploited by the repository's segment checker.

The supplied analyzer regenerates the existing certificate from the
pinned input files and checks the effects on both sides. It finds only
\textbf{32 distinct local effect signatures} among the 2,240 checkpoints.
The signatures consist of the tested register and the two sparse
increment vectors, with a separate halt signature.

These 32 signatures are not 32 complete proof states. Continuations
matter. Starting with the signature partition and repeatedly refining by
the two successor classes stabilizes at \textbf{1,423 classes}. The
analyzer checks explicitly that every refinement refines the preceding
partition and that the final partition preserves both local labels and
successor classes. It writes the full quotient certificate.

\begin{proposition}[Soundness of the checkpoint quotient]
Suppose two checkpoint nodes have the same halt/test label, the same
branch increment vectors, and successors in corresponding classes of a
graph congruence. When started with the same register file, their
checkpoint executions halt equivalently and preserve equality of the
register files at corresponding checkpoints.
\end{proposition}

\begin{proof}
If the label is halt, both terminate. Otherwise both test the same
register value and hence select the same branch. The decrement, if
performed, agrees; adding the same increment vector gives equal new
register files. The successor nodes again belong to the same class.
Induction on the number of checkpoint transitions transfers any finite
halting execution in either direction.
\end{proof}

\subsection{What this measurement does and does not buy}

The graph quotient has 817 fewer classes than the original certificate,
about 36.47\% fewer. The effect library has only 32 distinct entries. A
proof representation can exploit both facts, but these percentages are
\emph{not} measured reductions in Lean proof size or kernel time.
Every literal path still needs a justified binding to its checkpoint
node, including the root segment. A quotient without those bindings
would prove a statement about a detached graph, not about the source.

The recommended implementation has three layers: a verified conversion
from source paths to checkpoint effects, a compact graph-congruence
certificate, and the existing halting interface. Cache shared effect
proofs in named lemmas and store source addresses separately. Retain the
current checker as a comparison implementation until the exact K215
endpoint and all assumption reports pass through the new route.

The source analyzer and quotient certificate are finite Python artifacts,
not proof-assistant theorems. They supply a concrete target for that
refactor. They also prevent confusing the 179-node compiled path-program
DAG in the Lean source with the 202-subtree source representation used
by the machine construction; those are different representations.

\section{Completed searches for a 214-state machine}

\subsection{Pinned inputs and abstract contexts}

The search suite contains fourteen literal 215-state tables: primary
K215, CM215, commands120, composed, reencoded, zero-spare, zero-spare-z201,
zero-spare-transfer, MP-first transfer, both rho6 orders, zero-spare
MP-first, direct-transfer ST215, and rebalanced R41. Paths and hashes
are recorded in \path{reduction_search/inputs.json}. This is a selected
inventory, not a claim to enumerate every table or graph-isomorphism
class in the repository.

An abstract configuration consists of a control and the fifteen bits
at head offsets $-7,\ldots,7$. Writes and shifts update known cells
exactly; a newly exposed boundary bit may be either zero or one. Closing
these transitions gives an inductive overapproximation of reachable
windows. One initial domain uses every possible window at the start
control; the other uses the zero window, which includes the actual blank
tape but allows arbitrary distant exterior bits.

For changed rows, the filter seeks a positive common endpoint within
97 transitions on each side, without reading outside the known initial
window. At a nonterminal join it requires equal heads, equal whole known
windows, and related controls. Neither run has written outside the
window, so this gives whole-tape equality. When both prefixes halt,
terminal tape equality is unnecessary. A final move may leave the window;
the next read may not occur outside it before a join.

\subsection{Family A: inherited rows and one common branch}

For every ordered pair $(u,v)$ of distinct working controls, remove $v$
and redirect references to $v$ into $u$. All other rows retain their
mapped actions. On each symbol the keeper may use either of the two
mapped inherited actions. If exactly one branch agrees, the other branch
is also allowed every write, direction, and remaining target, including
the halt. There are 860 possible actions on such a free branch.

Every input has $215\cdot214=46,010$ ordered pairs. The total candidate
specification count for one initial domain is 8,854,464, and both domains
therefore produce \textbf{17,708,928 full-specification evaluations}.
No candidate survives. The count is of evaluated specifications; different
specifications can describe the same literal machine.

\begin{center}
\begin{tabular}{@{}lr@{}}
\toprule
Input & Specifications per domain\\
\midrule
Primary / commands120 / composed / rebalanced & 632,584 each\\
CM215 / zero-transfer / zero-MP-first & 630,872 each\\
Reencoded & 641,144\\
Zero-spare / z201 & 634,296 each\\
MP-first & 629,160\\
rho6-120 / rho6-210 & 636,008 each\\
Direct transfer & 620,600\\
\bottomrule
\end{tabular}
\end{center}

The Python auditor independently recomputes these combinatorial counts
from the tables; it does not merely accept the C++ search counters.

\subsection{Family B: independently synthesized branches}

For primary K215 and CM215, a second search permits both keeper actions
to change. Each branch is synthesized independently from all 860 actions.
Before rejoining, its new block may not read the opposite keeper branch.
This restriction makes independent synthesis sound as a filter but
excludes solutions that must alternate between the two new rows.

The search first tries zero-branch actions. If none passes for a pair,
there is no need to enumerate its one-branch actions. Any pair with
passing actions on both symbols would be sent to the full joint filter.
No such complete pair was found.

\begin{center}
\begin{tabular}{@{}llr@{}}
\toprule
Input & Initial domain & Individual branch evaluations\\
\midrule
K215 & All windows & 39,883,360\\
K215 & Zero window & 39,883,360\\
CM215 & All windows & 39,871,320\\
CM215 & Zero window & 39,871,320\\
\midrule
Total & & 159,509,360\\
\bottomrule
\end{tabular}
\end{center}

This total is emphatically \emph{not} a count of that many distinct
214-state machines. It counts individual branch specifications in the
restricted synthesis procedure.

\subsection{Controls and interpretation}

Two small synthetic tables provide nonvacuous controls. A table with
duplicate terminal roles produces ten Family-A survivors and 288
Family-B survivors; every survivor is freshly checked by an unchanged
independent repository checker. A table with incompatible halting and
forever-moving roles produces none. The packet retains all 32 large
search receipts and a fresh selected replay. A separate root review
recompiled the implementation, reran the controls, and reproduced the
primary all-window Family-A record exactly apart from elapsed time.

Failure of a bounded join does not show that a candidate has the wrong
blank behavior. The filter may encounter an abstract window that is
unreachable in the concrete execution, need a larger window or longer
join, or need a relation that changes the tape representation. Neither
search family covers coordinated predecessor rewrites, arbitrary source
recompilation, or every interacting two-row transformation. No
214-state impossibility theorem follows.

The practical lesson is to move beyond a single isolated merged control
with all other mapped rows fixed. Coordinated incoming-row edits,
source reassociation followed by allocation, or a tape marker that
retains the former role are better targets for the next search. A
candidate should carry the intended simulation relation from its
construction, so its correctness can use the small transport interface.

\section{Reproduction and integration}

The archive is organized as a repository-ready research packet. The
root \path{README.md} and \path{verify.py} describe exact commands and
distinguish inexpensive finite audits from optional compiler checks and
full search replay. The report source can be compiled with
\texttt{pdflatex}; no external bibliography service is needed.

\begin{center}
\begin{tabularx}{\textwidth}{@{}p{0.28\textwidth}X@{}}
\toprule
Path & Contents\\
\midrule
\path{report/} & This article as editable LaTeX and compiled PDF.\\
\path{formal/} & Generic Lean/Coq transport, concrete CM215 Lean facts,
invariant, simulation, and dense-machine bridge.\\
\path{machines/} & Exact K215 and CM215 dense tables and native references.\\
\path{source_analysis/} & Pinned source inputs, reference generator,
structural analyzer, and graph quotient certificate.\\
\path{reduction_search/} & C++ search, fourteen exact inputs, independent
controls, standard-library Python driver, and all bounded-search receipts.\\
\path{integration/} & Small endpoint wrapper for the existing ProveIt
project and integration instructions.\\
\path{evidence/} & New successful compiler logs, assumption reports,
fresh finite audits, and explicit scope records.\\
\bottomrule
\end{tabularx}
\end{center}

No compiled search executable is required or shipped. Kernel artifacts
can be regenerated with the stated versions. The generic Coq check and
the concrete Lean extension are separate checks. The original
185-module Lean and 237-source Rocq calibration records remain the
baseline evidence; this packet does not relabel them as freshly rerun.

Integration should preserve the existing K215 endpoint and add CM215
through a short theorem composition. A successful complete project
build should print assumptions for the resulting CM215 endpoint as
well as for the new transport theorem. That final project-level audit
will verify the imported calibration and new extension together.

\section{Further research questions}

\begin{enumerate}
\item \textbf{A verified local-certificate format.} Can the three-clause
CM215 proof be generated from a small finite-window certificate whose
soundness is shared by Lean and Rocq? A useful first target is a format
containing incoming-row predicates, whole-tape block equalities, and
positive progress, with exact table hashes outside and literal rows
inside the kernel boundary.

\item \textbf{Interacting keeper branches.} Can the second reduction
search be extended to permit alternating between both synthesized rows,
while solving their finite constraints jointly? This would remove a
specific stated restriction, rather than merely increasing numerical
search bounds.

\item \textbf{Coordinated marker rewrites.} Can the selector role reused
by CM215 share an existing physical control after a few predecessor
edits? The search relation may retain origin in a neighboring marker
instead of requiring the same tape at every one-step boundary. The
three-clause invariant suggests exactly which information must survive.

\item \textbf{A zero-offset source checker.} Parameterize decoding depth
and prove that the current K215 certificate is a zero-offset instance.
Measure generated source bytes, kernel object bytes, peak memory, and
compile time separately. All four are relevant proof costs.

\item \textbf{A quotient-based source certificate.} Formalize the 32
effect templates and 1,423-class checkpoint quotient, including every
literal-path binding and the root segment. Does the smaller semantic
graph actually shorten the checked source proof, or do address proofs
dominate the remaining cost?

\item \textbf{Proof-aware source optimization.} Search jointly for
state count and new proof burden. A table requiring six local row facts
may be a better formalization target than one with fewer primitive
occurrences but a new scheduler invariant. A useful optimizer should
report both a physical state count and a proposed reusable theorem
interface.

\item \textbf{Arithmetic internalization.} Isolate the elementary
machine transport and compiler layer in a weak arithmetic formalization,
then connect it to the separately established internal consistency
statement. Keep this milestone distinct from external Lean or Rocq
proofs of the semantic equivalence.
\end{enumerate}

The immediate deliverable is a reusable route to a second 215-state
endpoint. A 214-state construction remains open after these particular
searches, and the packet specifies which transformations and proof
templates have actually been tested.

\section*{Attribution}
\addcontentsline{toc}{section}{Attribution}

K215, CM215, the numerical-calculus calibration, source/compiler
architecture, original marker transformations, and existing verification
framework are inherited from ProveIt and the upstream authors credited
in its notices. The existing CM215 machine is not claimed as a new
construction here. This packet contributes the compact invariant-based
formal transport development, generic checked simulation lemmas,
source-certificate analysis, and the explicitly scoped search extension.
Copied source and table notices are retained. No global novelty claim is
made for the general simulation or partition-refinement principles.

\begin{thebibliography}{12}
\bibitem{verification} ProveIt authors.
\emph{Formal correctness of the exact 215-state machine}.
Pinned verification index, Lean/Rocq build and assumption records.
\href{\repo/Computability/BusyBeaver/Research/ZFC/k225-formal/verification/k215/README.md}{Repository source}.
\bibitem{phase3} ProveIt authors.
\emph{Shared lemmas and one enumerator theorem}, with phases 1--3.
\href{\repo/Computability/BusyBeaver/Research/ZFC/k225-formal/verification/k215/simplification/phase3/README.md}{Repository source}.
\bibitem{endpoint} ProveIt authors.
\emph{KEndpoint215.lean}.
\href{\repo/Computability/BusyBeaver/Lean/BusyBeaver/ZFC/KEndpoint215.lean}{Exact Lean endpoint}.
\bibitem{cm215} ProveIt authors and credited contributors.
\emph{CM215: conditional origin markers and ZFC consistency}.
\href{\repo/Computability/BusyBeaver/Research/ZFC/candidate-k215/alternatives/conditional-marker/README.md}{Machine, ordinary proof, and retained evidence}.
\bibitem{sourceproof} ProveIt authors.
\emph{KSource215Data.lean}, \emph{RMBisim.lean}, and
\emph{RMOffsetBisim.lean}.
\href{\repo/Computability/BusyBeaver/Lean/BusyBeaver/ZFC/KSource215Data.lean}{Source certificate};
\href{\repo/Computability/BusyBeaver/Lean/BusyBeaver/ZFC/RMOffsetBisim.lean}{Offset checker}.
\bibitem{rebalanced} ProveIt authors and credited contributors.
\emph{K215-R41: rebalanced source addresses}.
\href{\repo/Computability/BusyBeaver/Research/ZFC/candidate-k215/alternatives/rebalanced/README.md}{Delivered packet};
\href{\repo/Computability/BusyBeaver/Research/ZFC/candidate-k215/alternatives/rebalanced/INTAKE.md}{Intake status}.
\bibitem{directtransfer} ProveIt authors and credited contributors.
\emph{ST215: direct 35-leaf transfer through contextual reassociation}.
\href{\repo/Computability/BusyBeaver/Research/ZFC/candidate-k215/alternatives/direct-transfer/INTAKE.md}{Intake status}.
\bibitem{commands120} ProveIt authors and credited contributors.
\emph{A 215-state machine with DROP/NIP/MP command order}.
\href{\repo/Computability/BusyBeaver/Research/ZFC/candidate-k215/alternatives/commands120/README.md}{Repository source}.
\bibitem{zerospare} ProveIt authors and credited contributors.
\emph{Z215: zero-valued spare blocks and a two-truth startup}.
\href{\repo/Computability/BusyBeaver/Research/ZFC/candidate-k215/alternatives/zero-spare/INTAKE.md}{Intake status}.
\end{thebibliography}

\end{document}
